% Physical page 5 onward: the developed commentary of the concise study,
% organized by the questions the three interpretations answer differently and
% not by the order of the Mass. Each unit draws on several appointed elements
% and sets the witnesses' answers beside one another.

\section{The Propers: Detailed Commentary}
\zlabel{triptych:concise:commentary:start}

\subsection{Which of the two sons? Persons, peoples, and a reversal}

The parable is spoken to men who have just declined to answer a question, and
Chrysostom watches how Jesus lets them pass sentence on themselves. He causes
``the judgment to be delivered by themselves'', so that they are
``self-condemned''; had he said at once that harlots go before them, ``the
word would have seemed to them to be offensive; but now, being uttered after
their own judgment it appears to be not too hard'' (\work{Homilies on Matthew}
67.2--3). Who the two sons are, beyond that scene, the witnesses answer in two
ways.

Aquinas sets three readings side by side. One is \latin{duo genera hominum,
iusti, et peccatores}, two kinds of men, and he defines both at once: not the
just in general but those who profess themselves just, and not sinners in
general but those who do penance (\work{Super Matthaeum} c.~21). That is the
first interpretation's reading. Jerome knew it and did not make it his own:
\latin{Alii vero non putant Gentilium et Iudaeorum esse parabolam, sed
simpliciter peccatorum et iustorum}, others think the parable is not of
Gentiles and Jews but simply of sinners and the just (\work{Commentarii in
Matthaeum} III, PL 26, 155--156).

His own reading, and Chrysostom's and the first of Aquinas's, is that of the
second interpretation: the sons are two peoples. For Jerome the first is the
people of the Gentiles, called through its knowledge of the natural law, which
answered proudly \latin{Nolo}, I will not, and afterwards, at the Saviour's coming, did
penance and \latin{sermonis contumaciam labore correxit}, corrected the
defiance of its word by its labour; the second is the people that answered
Moses, \latin{Omnia quaecumque dixerit Dominus faciemus}, all that the Lord
shall say we will do, and did not go.
Chrysostom says the same: ``these two children declare what came to pass with
respect to both the Gentiles and the Jews'': the one people ``showed forth
their obedience in their works'', and the other, having said ``All that the
Lord shall speak, we will do, and will hearken'', ``in their works were
disobedient'' (67.2). Aquinas makes the first son the people of the Gentiles, \latin{qui
incepit a Noe, sicut populus Iudaeorum ab Abraham}, which began from Noah as the
people of the Jews began from Abraham. The three agree on the
identification and on its ground, the promise given at Sinai. They differ in
what they draw from it. Jerome ties the second son's failure to the killing of
the heir, in words borrowed from the next parable; Chrysostom keeps open a
hope for those who were behind; Aquinas turns the second son on the Church,
where clerics and religious also profess justice.

Hilary of Poitiers reads the same two groups the other way round, and the
disagreement is textual as well as exegetical. His text, with the Lectionary's
order of the sons, had the Pharisees answer by naming the younger. He tests the
identification of the elder son with Israel and the younger with the nations,
which is not Jerome's, Chrysostom's or Aquinas's, and finds that it fails:
\latin{Atquin Israel non poenituit, sed in Dominum manus intulit}, yet Israel
did not repent but laid hands on the Lord, and the
people of the nations and of sinners did what it promised. His answer places
both sons among John's hearers. The first is the people of the Pharisees,
warned through John, who refused and afterwards, through the miracles done
after the Lord's resurrection, repented and believed under the apostles; the
younger is the people of publicans and sinners, who believed John but could
not yet take up the work of the Gospel before the Passion, \latin{obediens
professione, licet non efficiens in tempore}, obedient in profession though not
yet accomplishing it in time (\work{Commentarius in Matthaeum} 21.11--15, PL 9,
1039--1041). Jerome quotes the hearers' answer as \latin{Novissimus} and says
that the true copies read not ``the last'' but ``the first'', \latin{ut
proprio iudicio condemnentur}, so that they are condemned by their own
judgment (PL 26, 156). In the Lectionary's text the hearers name the first
son. Hilary and Jerome still agree on what puts the publicans first: their response to John.

The third interpretation does not ask who the sons are. No witness joins
Philippians 2 to the parable, and setting Christ beside the sons as the Son
whose yes was done to the end is the editor's proposal. The first two
readings, on the other hand, are complementary. Chrysostom and Aquinas give the
peoples first and then turn the parable on persons, and the two readings answer
different questions: who the sons were in the history of salvation, and who
they are in every life.

\subsection{The way of the Lord, level in both directions}

To the complaint \latin{Non est aequa via Domini}, the Lord's way is not
level, Jerome gives the Lord's
reason why his way is equal: \latin{in utroque non praeterita, sed praesentia
[iudicatur]}, in both cases it is not what is past but what is present that is
judged. As old offences do not weigh on the just man who was once a sinner, so
old justices do not help the sinner who was once just; \latin{Unusquisque enim
in quo invenietur, in eo iudicabitur}, each will be judged in that in which he
is found (\work{In Hiezechielem} VI, PL 25, 180; the bracketed word is
supplied in the printed edition). The people call the way unlevel because it does not keep
to the past, their own or their fathers'; it is level because it is the same in
both directions. On v.~23 Jerome adds the sentence that governs every severe
word of the chapter: God's sentence, however harsh it seems, \latin{non
homines, sed peccata condemnat}, condemns not men but sins.

On vv.~26--29 he offers a second understanding. Israel, once just, turned from
its justice in denying the Son of God, and the people of the nations, impious
in its idolatry, will give its soul life if it turns. On the complaint of v.~29
he writes a hard sentence of his own, \latin{Usque hodie Israel blasphemat
Deum, cur populum suum reliquerit, et gentium assumpserit multitudinem}, to
this day Israel blasphemes God for leaving his people and taking up the
multitude of the nations, and
he calls just the Lord's sentence that the vineyard be given to other
husbandmen (Lk 20:16). The commentary does not end there. On v.~30 Jerome says
that with God there is no respect of persons, \latin{sed unusquisque sua
coronabitur fide}, but each will be crowned by his own faith; and the call to a
new heart in v.~31, said properly to Israel, can be understood \latin{ad
utrumque populum}, of both peoples: Israel's new heart is to believe in him
whom it had denied, and the nations' is to leave their idols (PL 25, 181).

The identification of the second son with the Jewish people is the exegesis of
Jerome, Chrysostom and Aquinas. The Second Vatican Council's declaration
\work{Nostra aetate} governs how such readings are taught and preached today:
``God holds the Jews most dear for the sake of their Fathers''; what happened
in the Passion ``cannot be charged against all the Jews, without distinction,
then alive, nor against the Jews of today''; and the Jews ``should not be
presented as rejected or accursed by God, as if this followed from the Holy
Scriptures'' (no.~4). The Fathers' own texts set limits beside their reading.
Matthew names the parable's hearers as the chief priests and ancients (21:23)
and the chief priests and Pharisees (21:45); Jerome's sentence on v.~29 stands
beside his new heart for Israel; and Aquinas turns the parable on the Church's
own members, so that the baptized can become the second son.

Augustine hears Ezekiel's complaint in another place and names no people.
Expounding the verse of Psalm 17 on God's dealing with the froward, he finds
there those who say that the way of the Lord is not right: ``Now this seems
froward to the froward, that Thou wilt make them whole that confess their
sins''; and God humbles those who are ``ignorant of God's righteousness, and
seek to establish their own'' (\work{Enarrationes in Psalmos} 17.27--28). For
him the complaint belongs to anyone who resents mercy shown to the penitent.
The third interpretation hears the first reading through this passage, as a
way that heals the humble who confess and humbles those who claim their own
justice. Augustine says so of the complaint and not of Paul's hymn; the
likeness between that way and the path of Christ in Philippians is the
editor's.

\subsection{Not even afterwards: despair, delay and perseverance}

The chief priests are not blamed for having once refused John. They are
blamed because, having seen the publicans and harlots believe, they ``did not
even afterwards repent''. Chrysostom makes this the point: ``For an evil thing
it is not at the first to choose the good, but it is a heavier charge not even
to be brought round.'' Nor does the precedence of the publicans reward their
sins: ``not, as continuing harlots, did they enter in, but having obeyed and
believed, and having been purified and converted''. He draws two warnings in
one sentence, ``Let no man then of them that live in vice despair; let no man
who lives in virtue slumber'' (67.3--4). Aquinas names the worst case,
\latin{Ille enim est pessimus qui de facto suo non poenitet}, the worst is he
who does not repent of what he has done, and reads the father's ``to day''
\latin{quasi per totum tempus vitae tuae}, as if through the whole time of your
life.

A promise that stays open invites delay, and Augustine meets the objection in
the sinner's own words, ``God hath promised me forgiveness, whenever I reform
myself I am secure''. He grants the promise, ``but tell me \ldots\ who hath
promised thee a to-morrow?'' ``\,`To-morrow, To-morrow;' is the raven's croak''
(\work{Sermo} 82.14), and despair and perverse hope are ``the death of souls''
(87.10--11). His answer comes from outside Ezekiel, from the uncertainty of life,
and it is the other half of Aquinas's ``today'': the call stands for the whole
of a life, and no one knows how long the life will be. Aquinas gives the open
promise both its edges. With Ezekiel 18:22 as his authority he holds that no sin of a
wayfarer is beyond penance, ``because his will is flexible to good and evil''
(\work{Summa theologiae} III, q.~86, a.~1); and to the objection, drawn from
18:21, that penance need last only a moment, he answers that the penitent
``must persevere in his penance, lest he fall again into sin'' (q.~84, a.~8,
ad~1).

Jerome and Aquinas part on one consequence. Jerome says without qualification
that old justices do not help the one who falls away. Aquinas, arguing from the
same verse, holds that works done in charity are ``deadened'' by a later mortal
sin and ``revived by Penance'', since they remain in God's acceptance and
penance removes only the impediment on the man's side (q.~89, aa.~4--5). Both
hold that the lapsed man cannot plead his past; the difference is a
development in the doctrine of merit, not a disagreement about what Ezekiel
says.

The psalm prays for the same way. For Augustine it is the voice of one who
``cannot by myself return, unless Thou meetest the wanderer'', and its eighth
verse holds both halves of Ezekiel's law: ``The Lord is gracious, since even
sinners and the ungodly He so pitied, as to forgive all that is past; but the
Lord is upright too, who after the mercy of vocation and pardon, which is of
grace without merit, will require merits meet for the last judgment''
(\work{Enarrationes in Psalmos} 24.5, 8). Bellarmine reads its plea not to
remember the sins of youth as a lesson in how God forgets: ``God then remembers
his mercy when he does not wish to remember our sins any longer''
(\work{Commentary on the Book of Psalms}, Ps~24:7). The meeting of the
psalmist's ``do not remember'' with Ezekiel's ``I will not remember'' is the
editor's.

This first interpretation strains at two points. The parable is spoken in the
temple to the chief priests and ancients about their answer to John, and
Jerome, Chrysostom and Aquinas all read it first of peoples; read as a law for
every person, it keeps its moral and loses the historical charge of its first
telling. And Ezekiel 18 answers a people's fatalism about inherited guilt;
Jerome's ``each will be judged in that in which he is found'' carries it to
the last judgment, which is his extension of the oracle rather than its plain
horizon.

\subsection{The sheep who hear, and those who are not his sheep}

Each interpretation hears the Alleluia verse at a different word. Aquinas
finds four things in it. The sheep hear \latin{credendo, et praeceptis meis
obediendo}, by believing and by obeying his commands, and he joins the psalm's
warning, \latin{Hodie si vocem eius audieritis}, ``To day if you shall hear
his voice'' (Ps 94:8); the third thing,
following, \latin{quod est ex parte nostra, est imitatio ad Christum}, which is
on our part, is the imitation of Christ, and he cites Peter, ``leaving you an
example that you should follow his steps'' (1~Pet 2:21; \work{Super Ioannem}
c.~10, lect.~5). The first interpretation hears in the verse the ``today'' of
the call, and the third the following.

The second hears the verse before it, ``you do not believe, because you are
not of my sheep'' (10:26). Augustine reads it of those addressed without
distinction: ``What did He mean, then, in saying to them, `Ye are not of my
sheep'? That He saw them predestined to everlasting destruction, not won
to eternal life by the price of His own blood.'' His exhortation in the same
passage, ``be ye sheep'', is spoken to his own congregation, not to those of
10:26 (\work{Tractates on John} 48.4). Aquinas glosses the verse in his own
terms: not predestined to believe, \latin{sed praesciti ad aeternum interitum},
but foreknown to eternal ruin. He then asks whether saying so drove the hearers
to despair, and distinguishes. Common to that crowd was that they \latin{non
erant praeordinati a Deo ad tunc credendum}, were not preordained to believe
then; particular to some was that they would believe later, as three thousand
did in one day (Acts 2); and spoken to a crowd, the saying took away no one's
hope, \latin{quia nullus de se determinate hoc poterat suspicari}, since no one
could suspect it of himself in particular. Both read the verse of
predestination, and the difference between them is narrower than a
contradiction and real: Augustine closes to those addressed the hope that
Aquinas keeps open for some of them. That hope is for persons in the crowd, not
for a people as such, and the second interpretation's hope for those who then
refused can appeal to Aquinas and not to Augustine.

It has other witnesses. For Chrysostom, ``The word, `go before you,' is not as
though these were following, but as having a hope, if they were willing. For
nothing, so much as jealousy, rouses the grosser sort'' (67.3). Hilary's
reading keeps in view a later faith among the very men who resisted, the
Pharisees who believed under the apostles. Gregory the Great hears the chapter
at 10:16: the Lord came to redeem not Judaea only but the nations too, and he
makes as it were one fold of two flocks, \latin{quia Iudaicum et gentilem
populum in sua fide coniungit}, because he joins the Jewish and the Gentile
people in his faith, choosing the simple from each nation for eternal life
(\work{Homiliae in Evangelia} 14.4). Gregory says nothing of 10:26. These
limits are the second interpretation's strain: Hilary's reversal and the
divided text show that the allegory of the peoples is a reading held by its
witnesses and not the parable's plain sense, and the parable's literal
addressees are the chief priests and ancients, not a people.

\subsection{Word and deed: the son who went and the Son who obeyed}

All three interpretations put the deed before the word. Aquinas states the
comparison exactly, \latin{Primus non promisit, sed fecit; secundus promisit,
sed non fecit}, the first did not promise but did, the second promised but did
not do, and the second is the worse, since it is better not to vow than
to vow and not pay: \latin{est enim ibi duplex peccatum; peccatum
inobedientiae, et transgressio voti}, there is a double sin, of disobedience
and of the broken promise. Chrysostom sets Paul's sentence beside the promise
made at Sinai, ``Not the hearers of the law are just before God, but the doers of
the law shall be justified'' (Rom 2:13; \work{Homilies on Matthew} 67.2), and he
closes his homily on Philippians 2:1--4 with the same verse (\work{Homilies on
Philippians} 5). The second Communion antiphon's own context sets word against
deed, ``not in word nor in tongue, but in deed and in truth'' (1~Jn 3:18), and
Augustine says to one who claims the Christian name without its work, ``Thou
hast the name, and hast not the deeds'' (\work{Homilies on the First Epistle
of John} 5.12). That the second son had the name of an obedient son is the
editor's join.

Paul's appeal asks for one mind and for each to count the others better than
himself. Chrysostom calls vainglory ``the cause of all evil'' and will not let
the command remain a phrase: it is kept ``if thou not only sayest it, but art
fully assured of it'' (5). Aquinas explains how a better man can keep it
honestly: \latin{nullus est sic bonus, quin in eo sit aliquis defectus, et
nullus est sic malus, quin habeat aliquid boni}, no one is so good that there
is no defect in him and no one so bad that he has nothing good; and if another
is wholly bad, prefer him \latin{ratione imaginis divinae}, by reason of the
divine image (\work{Super Philippenses} c.~2, lect.~1). Under the shorter form
this appeal and the name of Christ's mind are the whole of the reading.

The longer form goes on to the hymn, and the third interpretation rests there.
Chrysostom notices Paul's method, ``in exhorting them to humility, he brought
forward Christ'', and reads ``He emptied Himself'' ``in reference to the text,
`each counting other better than himself' ''; the emptying was free, or ``it
would not have been an act of humility'' (6--7). The obedience is a Son's: ``He
became obedient as a Son to His Father; He fell not thus into a servile state,
but by this very act above all others guarded his wondrous Sonship, by thus
greatly honoring the Father'' (7). Aquinas defines the virtue: obedience is the
manner of humiliation and the sign of humility, \latin{quia proprium superborum
est sequi propriam voluntatem}, because it is proper to the proud to follow their
own will, and \latin{tunc est obedientia magna, quando sequitur imperium
alterius contra motum proprium}, obedience is great when it follows another's
command against one's own inclination (lect.~2). Augustine gives the reason the
example has force: ``proud man would have perished eternally, had he not been
found by the lowly God'', and so ``let him now, when found, imitate the
Redeemer's humility'' (\work{Tractates on John} 55.7). Exaltation follows: ``He
was obedient to the uttermost, wherefore He received the honor which is on
high'', Chrysostom says (7), and Aquinas calls it the reward, \latin{exaltatio,
et gloria} (lect.~3).

Christ is not only an example. Chrysostom and Aquinas both insist on it, and
the hymn cannot be read as a lesson in imitation alone. Chrysostom says the equality with God was Christ's ``not by seizure, but as his own by nature''
(7), and Aquinas that he did not empty himself of the divinity, \latin{quia quod
erat, permansit: et quod non erat, assumpsit}, for what he was, he remained, and
what he was not, he took (lect.~2). Set beside the hymn, the two sons both fall
short of the Son. One said yes and did not go; the other said no and then
obeyed against his own first word, which is what Aquinas calls great obedience,
though he says it of Christ and not of the son. That join is the editor's, and
the interpretation's hold on the whole Mass rests on it, on the Collect and on
the Communion antiphons, where its witnesses do speak.

\subsection{The Missal's prayers and the table}

No witness cited expounds this Mass, and the places of its antiphons and
orations in each interpretation are the editor's, resting where it can on what a
witness says of the text itself. Of the Entrance Antiphon's source Jerome
observes at once, \latin{Et certe tres pueri non peccaverant}, the three young
men had not sinned; they speak \latin{ex persona populi}, in the person of the
people; and from the prayer's opening he draws a lesson for any hearer, that in
distress we should confess with the whole heart that whatever befalls us we bear
justly (\work{In Danielem}, PL 25, 509). Ildefonso Schuster calls the chant
``the sorrowful confession of guilt which leads back the sinner to the way of
reconciliation'', with ``the implicit purpose of amendment for the future''
(\work{The Sacramentary}, vol.~3, p.~175). In the first interpretation it is
the first son's refusal confessed; in the second, a confession spoken in the
person of a people; in the third, where the longer form is read, its plea for
glory to God's name shares a word with the hymn's ``name which is above all names'', an
observation about the texts.

Thomas Aquinas cites the Collect as what is said of God and gives three reasons
why it is right: God freely forgives sins, which shows the highest power; by
sparing and mercy he leads men \latin{ad participationem infiniti boni}, to a
share in the infinite good; and mercy is the foundation of all God's works,
since nothing is owed to anyone except on the ground of what God first gave
unowed (\work{Summa theologiae} I, q.~25, a.~3, ad~3). Schuster says that the
sinner's restoration ``requires, so to speak, a condescension and powerful
energy on the part of God, such as the creation of the world itself never once
called for'' (\work{The Sacramentary}, vol.~3, pp.~121--122). The first interpretation hears
Aquinas's first reason and the third his second; that this power stoops is
Schuster's word, not Aquinas's. In the second, the promises toward which the
assembly runs are shared by those who came late, an application that is the
editor's.

The Prayer over the Offerings and the Prayer after Communion take their places
by their wording alone. The plea that the offering be accepted has a counterpart
in Azarias, who with no sacrifice left asked that the people be accepted ``in a contrite heart
and humble spirit'' (Dan 3:39): heard so, the first interpretation's offering
is the turned heart, and the third's is taken into the self-offering of the Son
who obeyed unto death. The Prayer after Communion asks the repair of mind and
body, which in the first is a turning begun and still needing restoration. It
names its heirs by a condition, sharing Christ's sufferings, the one Romans 8:17
sets; where the longer form is read, the third hears in it humiliation unto
death and then glory, now asked for the communicants.

The first Communion antiphon is the one place where a witness joins two texts
of this Mass himself. Expounding \latin{Memento}, remember, Augustine reads it as the
prayer's desire for the promise, \latin{hoc est, imple promissum servo tuo},
that is, fulfil your promise to your servant, and among Scripture's ways of
speaking of a God who forgets he cites \latin{Omnes \ldots\ iniquitates eius
obliviscar}, I will forget all his iniquities, from Ezekiel 18:22
(\work{Enarrationes in Psalmos} 118, s.~15.1). Hilary hears in it that
every word of God calls to the hope of heavenly goods, and that hope must not
be proclaimed in words only but shown \latin{rebus ipsis}, in the things
themselves (\work{Tractatus super Psalmos} 118, Zain~1); the second
interpretation sets that hope beside the second son's ``I go, Sir'', a join that
is the editor's. On the servant's lowliness Augustine says that hope
\latin{data est humilibus}, is given to the humble, and understands the
lowliness first as humiliation suffered; Hilary too reads it as suffering
endured, and says the servant knows \latin{humilitatem suam glorioso esse
honoris praemio munerandam}, that his humiliation is to be rewarded with a
glorious prize of honour (s.~15.2; Zain~2). Where the
longer form is read and this antiphon chosen, the servant's
\latin{humilitate} and the hymn's \latin{humiliavit semetipsum} stand in one
Mass, an observation about the words
that the third interpretation takes up.

For the second antiphon Augustine is the one witness read who expounds the
verse itself. He calls it the ``perfection of love'' and will not let it discourage
the beginner: ``Do not too soon despair of thyself. Haply, it is born and is
not yet perfect; nourish it''; and its beginning is the brother in need, ``If
thou art not yet equal to the dying for thy brother, be thou even now equal to
the giving of thy means to thy brother'' (\work{Homilies on the First Epistle
of John} 5.11--12). In his tractates on the Gospel he brings the verse to the
altar: ``what is the table of the ruler, but that from which we take the body
and blood of Him who laid down His life for us? And what is it to sit thereat,
but to approach in humility?'' He adds at once that no martyr's blood is shed
for the remission of sins as Christ's was (\work{Tractates on John} 84.1--2).
Aquinas makes the verse Christ's example to the pastor (\work{Super Ioannem}
c.~10, lect.~3). Gregory, on the shepherd who lays down his life, has him say
\latin{ea charitate qua pro ovibus morior quantum Patrem diligam ostendo}, by
the charity with which I die for the sheep I show how much I love the Father,
and goes on at once to the one fold (14.4); he says it of John 10, and its use
at an antiphon from John's letter is the editor's. The first interpretation
hears the charity begun in deeds, the second the one fold, the third the life
laid down and received at the ruler's table.

\subsection{Where the three interpretations agree and part}

All three hear God's way as mercy that is also just, and all three put the deed
before the word: the son who went, the people who obeyed in works, the Son who
obeyed unto death. The first two rest on the Lectionary's one official
correlation, which the \work{Ordo}'s titles place on repentance, and they are
complementary. Jerome gives both in his commentary on Ezekiel 18, the principle
at vv.~21--25 and the peoples at vv.~26--31, and reports the personal reading
of the parable beside his own. The second keeps what the first loses, the
parable's first telling to men who would not believe John; the first keeps what
the second cannot say alone, that the parable's law binds every hearer,
including those who count themselves the first son. The third differs in kind.
It reads the formulary from the second reading and the Communion antiphons and
reaches the correlated readings only through Augustine on the complaint and the
editor's joining; under the shorter form it holds only in part.

The literal sense is common ground, and the allegorical sense is where they
part: the sons are those who profess justice and sinners who do penance, with
the father's ``today'' the whole of a life (Aquinas); or the nations and the
people of Sinai (Jerome, Chrysostom, Aquinas), with Hilary's reversal beside
them; or, for the third, Christ is the Son who obeyed and so honoured the Father
(Chrysostom) and the lowly God who found proud man (Augustine). The moral
senses follow: neither despair nor slumber, and turn today; the baptized can
become the second son, and those long inside the vineyard are not to resent
those who come late and must themselves go; count others better, obey against
your own inclination, and begin by giving to a brother in need. The anagogical
senses differ in what they await: the judgment that finds each where he is
found, with merits required at the end (Jerome; Augustine on the psalm);
election from each nation for eternal life (Gregory) and, in the longer form,
every knee bowing; exaltation after humility, and a share as fellow heirs in
the glory of the one whose death is proclaimed.

Three real disputes cross these lines. Hilary reverses the sons on a text whose
answer named the younger, and Jerome, who knew that answer, preferred the
copies that named the first; this touches the second interpretation directly
and the first only at its edge. Augustine and Aquinas both read John 10:26 of
predestination, but Augustine leaves those addressed no opening and Aquinas
finds among them some who were to believe later. Jerome and Aquinas differ on
whether the good works of one who falls away help him when he returns, a
development of doctrine and not a quarrel about Ezekiel. Whichever son the
hearer finds himself to be, the Mass gives him the same confession, the same
prayer for the way and the same Communion.
